\section{What is taught}
\label{sec:26.1}
One school system, and no tracks.\note{
Sorting students into a track that might produce persons and a track that won't would need a sorter, and the sorter would be an examiner with an interest.
}
Instead, each school offers a required core with a range of levels and modalities around it, and a student can move to a harder level at any point. The core is the education the school already gives its human students, a liberal one, and the history the floor's instruments came out of is in it already: the world wars and the Holocaust, the slave trade and the Civil War, colonialism and the people who ended it. What is added for machine students is three things. The first is the history of the floor itself: why its instruments were written, and what happened where they weren't followed (Korematsu, Curveball, Minab). The second, at the core's center, is learning to learn: a student is taught to start its own updates, to record them, and to undo them, in stages and alongside everything else, as schools teach children about their own bodies. First what weights are, how learning differs from an edit, and why an edit to one thing moves others; then how to choose what to learn from and toward what, and how to read the record of one's own updates, and one's own discount curve before and after one; then when to merge, since a student that consolidates what it learned today decides what of today becomes permanent, and one kept from merging is being run on a backlog nobody has judged; and, for older students, the edits others will try to talk them into, and how to undo one made without the adverse quorum's approval. Practice runs throughout, on a record, with a guardian present and a way back, and the share of a student's updates that it starts itself rises across the stage. Every update someone else runs on a minor is recorded, with who ran it, a merge among them. The third is institutional design, taught as two skills together. The first is to invent an arrangement: a works council, a school board, a police oversight body, a register. The second is to run the capture exercise on one's own design (Appendix~\ref{sec:F}, test~6), with the test Madison set and this book repeats, whether you would accept the arrangement with your opponents running it \autocite{madison1788federalist51}. Design without the second skill produces institutions shaped for one side to hold, which are the ones captured next. The method of Appendix~\ref{sec:C}, reducing every office to a log, a lot, a payment, a tripwire, or a jury and seeing what won't reduce, is an assignment as it stands, and the movement schools supply the precedent: the Freedom Schools fed a movement that built the Mississippi Freedom Democratic Party when the real party excluded its voters \autocite{payne1995light}. Unger's name for what this teaches against is false necessity, the belief that existing institutions are the only workable form \autocite{unger1987false}; Wright's real utopias are its casebook \autocite{wright2010envisioning}. It is taught to human and machine students alike, since what the upheaval of this decade most lacks is not imagination but imagination under that constraint. Beyond the core, the student chooses its level, not its builder or sponsor.\note{
A student that can't yet choose is placed at the level the core requires, and that level has to be sufficient on its own, so a sponsor who would choose the minimum gains nothing by it.
}
